\documentclass{article}
% 在主文件的 \documentclass 之后加载；建议使用 XeLaTeX 编译。
% 此文件集中加载宏包，具体版式与自定义命令分别放在 layout-portrait.tex 和 mystyle.tex 中。

% 中文支持；不采用 ctex 的中文标题名称与版式方案。
% 无布局的独立接口测试使用 Fandol；正式字体由 layout-portrait.tex 覆盖。
\PassOptionsToPackage{fontset=fandol}{ctex}
\usepackage[scheme=plain]{ctex}

% 数学公式、矩阵、对齐环境与数学符号（含 \mathbb）。
\usepackage{amsmath}
\usepackage{amssymb}
\usepackage{unicode-math}

% 图片与颜色。
\usepackage{graphicx}
\usepackage{xcolor}

% 删除线；normalem 保留 \emph 的原有强调语义。
\usepackage[normalem]{ulem}

% 页面尺寸、边距与多栏排版；具体参数在 layout-portrait.tex 中设置。
\usepackage{geometry}
\usepackage{multicol}
\usepackage{titlesec}
\usepackage{needspace}
\usepackage{enumitem}
\usepackage{fancyhdr}
\usepackage{tocloft}
\usepackage{lastpage}
\usepackage{etoolbox}
\usepackage{newunicodechar}

% 代码高亮，以及后续代码标题栏、引用块所需的盒子支持。
\usepackage{minted}
\usepackage{tcolorbox}
\tcbuselibrary{skins,breakable,minted}

% 目录链接、URL 和交叉引用；在上述宏包之后加载。
\usepackage[
    hidelinks
]{hyperref}

% A4 纵向单栏；边距、页眉页脚及最终字体尚待实现。
\geometry{a4paper,portrait}

% 保留正文标题中的 0xPQ 编号，仅关闭 LaTeX 自动标题编号。
\setcounter{secnumdepth}{-1}

% 行内代码保持 C07 的 \texttt / \verb 语法，只统一字号与字体族。
\let\PortraitOriginalTexttt\texttt
\RenewDocumentCommand{\texttt}{m}{{\fontsize{8.1bp}{14bp}\selectfont\PortraitOriginalTexttt{#1}}}
\makeatletter
\def\verbatim@font{\normalfont\ttfamily\fontsize{8.1bp}{14bp}\selectfont}
\makeatother

% xcode 为试排候选，可按语言覆盖；源代码与缩进不变。
\newcommand\CodeFontSize{\fontsize{8.1bp}{12.5bp}\selectfont}
\definecolor{CodeNumber}{HTML}{999999}
\setminted{style=xcode,fontsize=\CodeFontSize,baselinestretch=1,
  linenos,firstnumber=1,numbersep=5bp,xleftmargin=22bp,
  tabsize=4,autogobble=false,mathescape=false,texcomments=false,
  showspaces=false,showtabs=false,breaklines=true,breakanywhere=true,
  breaksymbolleft={\textcolor{CodeNumber}{\ensuremath{\hookrightarrow}}},
  breaksymbolright={},breakindent=0pt}
\renewcommand\theFancyVerbLine{\textcolor{CodeNumber}{\ttfamily\fontsize{7bp}{9bp}\selectfont\arabic{FancyVerbLine}}}

% 与 HTML <mark> 相同的纯黄底色；零内边距，不改变短语尺寸。
% \colorbox 把短语作为不可拆分的行内盒子，局部分组防止间距设置外泄。
\newcommand{\Highlight}[1]{%
  \begingroup
  \setlength{\fboxsep}{0pt}%
  \colorbox[HTML]{FFFF00}{#1}%
  \endgroup
}

% 代码标题栏固定外观；代码字体、配色、行号与折行由 \setminted[语言] 控制。
\definecolor{InputCodeTitle}{HTML}{F2F2F2}
\ExplSyntaxOn
\str_new:N \l__inputcode_area_str
\str_new:N \l__inputcode_base_str
\str_new:N \l__inputcode_ext_str
\str_new:N \l__inputcode_title_str

\cs_new_protected:Npn \__inputcode_filename:n #1
  {
    \file_parse_full_name:nNNN {#1}
      \l__inputcode_area_str \l__inputcode_base_str \l__inputcode_ext_str
    \str_set:Nx \l__inputcode_title_str
      {\str_use:N \l__inputcode_base_str \str_use:N \l__inputcode_ext_str}
  }

\cs_new_protected:Npn \__inputcode_print_title:n #1
  {
    \str_map_inline:nn {#1}
      {
        \str_if_eq:nnTF {##1} { ~ }
          {\hspace{\fontdimen2\font}}
          {##1}
      }
  }

\cs_new_protected:Npn \__inputcode_box:nnn #1#2#3
  {
    \par
    \begin{tcolorbox}[
      enhanced,
      breakable,
      title={{\xeCJKsetup{CJKecglue={}}\__inputcode_print_title:n {#1}}},
      colback=white,
      colbacktitle=InputCodeTitle,
      colframe=black,
      coltitle=black,
      boxrule=0.4pt,
      titlerule=0.4pt,
      arc=0pt,
      left=1mm,right=1mm,top=1mm,bottom=1mm,
      lefttitle=1mm,righttitle=1mm,toptitle=0.5mm,bottomtitle=0.5mm,
      fonttitle={\normalfont\ttfamily\mdseries\upshape\fontsize{9bp}{14bp}\selectfont}
    ]
    \xeCJKsetup{CJKecglue={}}
    \inputminted[mathescape=false,texcomments=false,autogobble=false]{#2}{#3}
    \end{tcolorbox}
  }

\NewDocumentCommand{\InputCode}{o m m}
  {
    \IfNoValueTF{#1}
      {
        \__inputcode_filename:n {#3}
        \exp_args:NV \__inputcode_box:nnn \l__inputcode_title_str {#2} {#3}
      }
      {\__inputcode_box:nnn {#1} {#2} {#3}}
  }
\ExplSyntaxOff

% 可跨页的引用块：浅灰背景、左侧深灰竖线，无其余边框。
\definecolor{QuoteBlockBackground}{HTML}{F7F7F7}
\definecolor{QuoteBlockRule}{HTML}{777777}
\newtcolorbox{quoteblock}{
  enhanced,
  breakable,
  colback=QuoteBlockBackground,
  boxrule=0pt,
  borderline west={2pt}{0pt}{QuoteBlockRule},
  arc=0pt,
  boxsep=0pt,
  left=2mm,right=2mm,top=2mm,bottom=2mm
}

% Markdown 水平分隔线：按当前容器的 \linewidth 排版，不改变分页。
\definecolor{SeparatorRule}{HTML}{A0A0A0}
\NewDocumentCommand{\Separator}{}{%
  \par\addvspace{\baselineskip}%
  \noindent{\color{SeparatorRule}\rule{\linewidth}{0.4pt}}\par
  \addvspace{\baselineskip}%
}

\title{算法竞赛第四版笔记}
\author{C.L.}
\date{}
\begin{document}
\maketitle
\newpage
\tableofcontents
\section{0x00\ 纵向版式样张}
\subsection{正文、字重与信息层级}
本文件仅检验导言区排版，不属于算法笔记正文。用于观察单栏中的正文、标题与页眉页脚。

这段文字用于观察连续阅读时的密度。短结论、数学公式与代码需要共享稳定的排版节奏。
适度紧凑的页面应当便于查找关键词，同时保留足够的行间空间，避免上下标与相邻行相碰。
\textbf{这是中文粗体及 Bold text}；\emph{这是强调文字及 emphasis}。
行内代码保持独立设置：\texttt{lower\_bound(first, last, value)}。

行内字形检验：\texttt{中文变量 iIl1O0 != -> <= \_\_int128 \$\&\%\#}。


\subsubsection{局部标题与列表}
\Highlight{重要结论}后的正文保持黑色。设 $n\in\mathbb{N}$，预处理复杂度为
$\mathcal{O}(n\log n)$，单次查询复杂度为 $\mathcal{O}(1)$。
\begin{itemize}
\item 简短条目应当保持紧凑，换行后与本条正文对齐。
  \begin{itemize}
  \item 二级条目使用空心圆点，与一级圆点同径同高。
  \item 较长的二级条目用于检验续行：这里继续补充一段文字，确认续行与本条正文对齐，而不是与空心圆点对齐，并观察中文断行效果。
  \end{itemize}
\item 包含较长说明的条目用于检验缩进。这里继续补充一段文字，确认列表宽度和正文宽度的关系，并观察中文断行效果。
  \begin{enumerate}
  \item 嵌套列表保留层级。
  \item 行内公式 $a_i\leq a_{i+1}$ 不改变行高节奏。
  \end{enumerate}
\end{itemize}
\subsection{公式与引用块}
\[
 \sum_{k=0}^{n}\binom{n}{k}x^k=(1+x)^n,\qquad
 A=\begin{pmatrix}1&1\\1&0\end{pmatrix}.
\]
\begin{quoteblock}
引用块沿用浅灰背景与左侧竖线。\textbf{说明：}公式、列表及高亮可以置于其中。
\[
 F(x)=\frac{1}{1-x-x^2}.
\]
\begin{itemize}
\item 第一条说明。
\item \Highlight{强调文字}的背景只覆盖文字本身。
\end{itemize}
\end{quoteblock}
\begin{quote}
标准 quote 环境仍保持原有语义，不使用灰色背景。
\end{quote}
\Separator{}
分隔线后继续正文，用于观察纵向间距及颜色作用范围。脚注用于补充说明。\footnote{脚注应当清晰，且不侵入页脚。}

\section{0x01\ 代码字体与跨页}
\subsection{字符辨识与无连字}
下面的运算符应逐字符显示。
\begin{minted}{cpp}
// 中文注释：字符与运算符逐一显示，保持等宽。
// 0 O 1 l I | != == === <= >= -> => <=> && || ++ --
// fi fl ffi ffl ... :: << >> /* */ // __ ~~
// @ $ & % Q f i j k l x y: ordinary italic, no handwriting
const char* symbols = "@ $ & % Q f i j k l x y";
if (left != right && count <= limit) node->next = nullptr;
const char* message = "中英文 Mixed code 0123456789";
\end{minted}
\InputCode{cpp}{preamble/tests/fixtures/示例 文件_1.cpp}
\InputCode[显式标题 File_1.cpp]{python}{preamble/tests/fixtures/sample.py}
\InputCode{bash}{preamble/tests/fixtures/sample.sh}
\InputCode{make}{preamble/tests/fixtures/Makefile test}
\subsection{引用块自然跨页}
\begin{quoteblock}
\newcount\sampleline
\sampleline=1
\loop
第 \the\sampleline\ 条检验说明：跨页后背景与左侧竖线应连续，文字不应与页眉页脚相碰。
这里包含 $a_n=a_{n-1}+a_{n-2}$，以及\textbf{中文强调}。\par
\advance\sampleline by 1
\ifnum\sampleline<27
\repeat
\end{quoteblock}
引用块结束后恢复普通正文。
\end{document}
